---
format: yohaku/0.2
target: text
macro: '[[macros/raw]]'
---

## 正文
````text
"""Keep the X11 input service independent of Yohaku and tool process groups."""
import ctypes
import ctypes.util
import os
from pathlib import Path
import select
import shutil
import subprocess
import time

UNIT = 'yohaku-fcitx-x11.service'
MARKER = '# Managed by Yohaku X11 IME lifecycle support\n'


def session_environment(env):
    result = dict(env)
    runtime = Path(env.get('XDG_RUNTIME_DIR', f'/run/user/{os.getuid()}'))
    result['XDG_RUNTIME_DIR'] = str(runtime)
    if not result.get('DBUS_SESSION_BUS_ADDRESS'):
        if not (runtime/'bus').is_socket():
            raise RuntimeError('X11 IME requires the existing user session bus; no bus was started.')
        result['DBUS_SESSION_BUS_ADDRESS'] = 'unix:path='+str(runtime/'bus')
    return result


def command(arguments, env):
    result = subprocess.run(arguments, env=env, capture_output=True, text=True, timeout=20)
    if result.returncode:
        raise RuntimeError(f'IME service operation failed ({Path(arguments[0]).name}, exit={result.returncode}); inspect journalctl --user -u {UNIT}')
    return result.stdout.strip()


def bus_owner(env):
    result = command(['busctl','--user','call','org.freedesktop.DBus','/org/freedesktop/DBus',
                      'org.freedesktop.DBus','NameHasOwner','s','org.fcitx.Fcitx5'], env)
    return result == 'b true'


def unit_text(environment_file, executable):
    return MARKER+f'''[Unit]
Description=Yohaku X11 Fcitx service (independent of application lifetime)
StartLimitIntervalSec=60
StartLimitBurst=3

[Service]
Type=dbus
BusName=org.fcitx.Fcitx5
EnvironmentFile={environment_file}
UnsetEnvironment=WAYLAND_DISPLAY WAYLAND_SOCKET
ExecStart={executable} -D --disable=wayland,waylandim
Restart=always
RestartSec=2
TimeoutStartSec=15
TimeoutStopSec=5

[Install]
WantedBy=default.target
'''


def environment_text(env):
    rows = []
    for name in ['DISPLAY','XAUTHORITY','XDG_RUNTIME_DIR','DBUS_SESSION_BUS_ADDRESS']:
        if name not in env:
            continue
        value = env[name]
        if any(character in value for character in '\r\n\0'):
            raise ValueError('Invalid X11 session environment value')
        rows.append(name+'="'+value.replace('\\','\\\\').replace('"','\\"')+'"')
    return '\n'.join(rows)+'\n'


def install(env):
    config = Path(env.get('XDG_CONFIG_HOME', str(Path(env['HOME'])/'.config')))
    directory = config/'systemd/user'
    directory.mkdir(parents=True, exist_ok=True)
    environment_file = config/'yohaku/ime-x11.env'
    environment_file.parent.mkdir(parents=True, exist_ok=True)
    executable = shutil.which('fcitx5', path=env.get('PATH'))
    if not executable:
        raise RuntimeError('Fcitx5 is not installed; Yohaku did not install anything.')
    path = directory/UNIT
    if path.exists() and not path.read_text().startswith(MARKER):
        raise RuntimeError('Existing IME unit is not Yohaku-managed; refusing to overwrite it.')
    text = unit_text(str(environment_file), executable)
    changed = not path.exists() or path.read_text() != text
    environment_file.write_text(environment_text(env))
    environment_file.chmod(0o600)
    if changed:
        path.write_text(text)
        command(['systemctl','--user','daemon-reload'], env)
    command(['systemctl','--user','enable',UNIT], env)
    return path


class XimSelection:
    def __init__(self, display):
        self.library = ctypes.CDLL(ctypes.util.find_library('X11'))
        for name, arguments, result in [
            ('XOpenDisplay',[ctypes.c_char_p],ctypes.c_void_p),
            ('XDefaultRootWindow',[ctypes.c_void_p],ctypes.c_ulong),
            ('XInternAtom',[ctypes.c_void_p,ctypes.c_char_p,ctypes.c_int],ctypes.c_ulong),
            ('XGetSelectionOwner',[ctypes.c_void_p,ctypes.c_ulong],ctypes.c_ulong),
            ('XSelectInput',[ctypes.c_void_p,ctypes.c_ulong,ctypes.c_long],ctypes.c_int),
            ('XPending',[ctypes.c_void_p],ctypes.c_int),
            ('XNextEvent',[ctypes.c_void_p,ctypes.c_void_p],ctypes.c_int),
            ('XConnectionNumber',[ctypes.c_void_p],ctypes.c_int),
            ('XCloseDisplay',[ctypes.c_void_p],ctypes.c_int),
        ]:
            function = getattr(self.library,name)
            function.argtypes, function.restype = arguments, result
        self.connection = self.library.XOpenDisplay(display.encode())
        if not self.connection:
            raise RuntimeError('Cannot connect to the configured X11 display for IME readiness.')
        self.atom = self.library.XInternAtom(self.connection,b'@server=fcitx',0)
        self.library.XSelectInput(self.connection,self.library.XDefaultRootWindow(self.connection),1 << 22)

    def owner(self):
        return self.library.XGetSelectionOwner(self.connection,self.atom)

    def wait(self, timeout):
        deadline = time.monotonic()+timeout
        event = (ctypes.c_long*24)()
        while True:
            if self.owner():
                return
            remaining = deadline-time.monotonic()
            if remaining <= 0:
                raise RuntimeError('Fcitx D-Bus service started but XIM selection did not become ready; no client success is claimed.')
            while self.library.XPending(self.connection):
                self.library.XNextEvent(self.connection,ctypes.byref(event))
            if self.owner():
                return
            select.select([self.library.XConnectionNumber(self.connection)],[],[],remaining)

    def close(self):
        self.library.XCloseDisplay(self.connection)


def ensure_x11_ime(env):
    session = session_environment(env)
    selection = XimSelection(session['DISPLAY'])
    try:
        owns_bus, owns_xim = bus_owner(session), bool(selection.owner())
        if owns_bus and owns_xim:
            return 'existing service reused (not restarted)'
        if owns_bus and not owns_xim:
            service_pid = command(['systemctl','--user','show',UNIT,'-p','MainPID','--value'], session)
            bus_pid = command(['busctl','--user','call','org.freedesktop.DBus','/org/freedesktop/DBus',
                               'org.freedesktop.DBus','GetConnectionUnixProcessID','s','org.fcitx.Fcitx5'], session)
            if service_pid.isdigit() and int(service_pid) > 0 and bus_pid == 'u '+service_pid:
                selection.wait(8)
                return f'{UNIT} ready after pending XIM initialization'
        if owns_bus or owns_xim:
            raise RuntimeError('Existing Fcitx has only D-Bus or XIM on this display; refusing to replace it. Inspect session/display connection.')
        install(session)
        command(['systemctl','--user','start',UNIT], session)
        selection.wait(8)
        if not bus_owner(session):
            raise RuntimeError('Fcitx lost its D-Bus name while XIM was starting.')
        return f'{UNIT} ready (D-Bus + XIM); remains running when Yohaku exits'
    finally:
        selection.close()


if __name__ == '__main__':
    print(ensure_x11_ime(os.environ))
````

## 输出与关联
```yaml
output: "research/workflow-runtime/agent/yohaku_ime.py"
```
